\section{Introduction}
\label{sec:intro}

The exponential growth of task-specific deep learning models has made multi-task model serving a cornerstone of modern machine learning infrastructure. A highly attractive paradigm is Parameter-Efficient Fine-Tuning (PEFT) using Low-Rank Adaptation (LoRA) \cite{hu2021lora}, where a shared pre-trained backbone is augmented with lightweight, task-specific expert adapters. At inference time, rather than hosting separate full-parameter models or running massive Mixture-of-Experts (MoE) architectures with high latency \cite{shazeer2017moe, fedus2021switch}, models can be consolidated in weight-space via \emph{model merging} \cite{wortsman2022model, ilharco2022editing, yadav2023ties} or dynamically blended in activation-space via \emph{dynamic ensembling} \cite{sable2025, regcalmerge2025}. Test-time activation blending is particularly appealing as it runs expert paths in parallel and scales their outputs dynamically, matching the flexibility of MoE without requiring joint multi-task training.

Despite its flexibility, test-time dynamic model ensembling faces a critical, unresolved challenge: the target inference stream is often highly non-stationary, and the calibration data available at serve-time is extremely scarce. Typically, a local router must adapt its routing coefficients using a transductive calibration split of only $N = 16$ or $N = 64$ samples per task. Under this severe data scarcity, standard unregularized Empirical Risk Minimization (ERM) easily overfits to transductive noise. The learned temperature parameters diverge to extreme scales, resulting in overconfident routing policies that dispatch queries to inappropriate experts and suffer from catastrophic generalization collapse on the out-of-sample target stream.

To prevent such overfitting, recent approaches have attempted to regularize the test-time routing parameters. For example, Subspace-Anchor Boundary Learning (SABLE) \cite{sable2025} relies on static, hand-tuned temperature parameters, which fail to adapt to varying noise and task characteristics. To provide generalization guarantees, PAC-ZCA \cite{paczca2026} optimizes Gaussian prior and posterior distributions over unconstrained log-temperature parameters. However, we argue that this formulation is fundamentally mismatched with the geometry of ensembling. Ensembling coefficients are mathematically constrained to lie on the probability simplex $\Delta^{K-1}$. Modeling unconstrained log-temperatures with Gaussian distributions ignores these simplex boundaries. This misalignment leads to a "log-temperature explosion" or sudden entropy collapse, causing severe numerical instability and loose generalization bounds.

\begin{figure*}[t]
  \centering
  \begin{subfigure}[b]{0.48\textwidth}
    \centering
    \includegraphics[width=\textwidth]{fig1.png}
    \caption{Joint Mean Accuracy vs. Entanglement Sweep ($\rho$)}
    \label{fig:sweep}
  \end{subfigure}
  \hfill
  \begin{subfigure}[b]{0.48\textwidth}
    \centering
    \includegraphics[width=\textwidth]{fig2.png}
    \caption{Bar Comparison ($\rho = 0.0$ vs. $\rho = 0.33$)}
    \label{fig:bar}
  \end{subfigure}
  \caption{Empirical evaluation of Dirichlet-PAC and baselines. (a) Joint mean accuracy across an entanglement sweep $\rho \in [0.0, 0.5]$ shows that Dirichlet-PAC degrades gracefully and maintains high performance, while unregularized ERM and PAC-ZCA collapse. (b) Bar comparison on orthogonal ($\rho=0.0$) and overlapping ($\rho=0.33$) manifolds.}
  \label{fig:results_main}
\end{figure*}

In this work, we propose that empirical success in deep representation learning can be enhanced by building upon rigorous mathematical foundations and formal statistical guarantees. To this end, we reformulate test-time model ensembling as a simplex-constrained statistical learning problem and introduce \textbf{Dirichlet-PAC}. Instead of unconstrained Gaussian distributions, we model the ensembling weight vector $\boldsymbol{\alpha}_b \in \Delta^{K-1}$ directly as a random variable distributed according to a Dirichlet posterior. By minimizing a rigorous PAC-Bayesian bound formulated under McAllester's theorem \cite{mcallester1999pac, mcallester2003simplified}, we constrain the learned task-specific temperatures using the exact analytical Kullback-Leibler (KL) divergence between Dirichlet distributions over the simplex itself.

To drive this routing policy without requiring ground-truth labels, we introduce **Subspace Energy Projection (SEP)**. SEP is an unsupervised, task-agnostic dimensionality reduction protocol that extracts orthonormal bases for task feature manifolds via Singular Value Decomposition (SVD) on early-layer calibration activations. The projection energies of incoming queries onto these subspaces map directly to the concentration parameters of our Dirichlet routing policy.

Minimizing our Dirichlet PAC-Bayesian objective establishes a formal complexity-performance duality. The Dirichlet KL penalty naturally suppresses transductive noise, acts as an elegant routing-entropy regularizer, and prevents temperature divergence, guaranteeing smooth, cooperative activation blending on task boundaries.

We evaluate Dirichlet-PAC on a 14-layer Analytical Coordinate Sandbox (ICS) across 10 random seeds. Our contributions are summarized as follows:
\begin{itemize}
    \item \textbf{Simplex-Constrained PAC-Bayesian Theory:} We present the first mathematically rigorous learning-theoretic framework for test-time model ensembling that operates directly on the probability simplex $\Delta^{K-1}$ using a Dirichlet policy.
    \item \textbf{Analytical Dirichlet KL Complexity Control:} We derive and present the exact analytical Dirichlet KL divergence within the PAC-Bayesian bound, providing a closed-form complexity penalty that stabilizes log-temperature optimization and prevents generalization collapse.
    \item \textbf{Unsupervised Subspace Projection with Energy Normalization:} We introduce Subspace Energy Projection (SEP) combined with a novel energy-normalization protocol, which projects activations onto SVD subspaces and maps task coordinates to the simplex.
    \item \textbf{Superior and Stable Serving Performance:} Under extreme data-scarcity ($N=64$), Dirichlet-PAC achieves an average accuracy of \textbf{77.88\%} on orthogonal and \textbf{76.32\%} on overlapping task manifolds. Among learned routers, Dirichlet-PAC significantly outperforms its uncalibrated prior SABLE (SEP-Block) Norm by \textbf{+4.56\%} (orthogonal) and \textbf{+4.80\%} (overlapping), and outperforms standard unregularized Temp-Only ERM by \textbf{+1.76\%} (orthogonal) and \textbf{+0.65\%} (overlapping), while dramatically reducing optimization variance and resolving representation corruption under high representation noise.
\end{itemize}
